\documentclass[11pt]{article}
\usepackage[margin=1in]{geometry}
\usepackage{amsmath,amssymb,amsthm,mathtools}
\usepackage{booktabs}
\usepackage{graphicx}
\usepackage{xcolor}
\usepackage{microtype}
\usepackage[hidelinks]{hyperref}
\usepackage{enumitem}

\newtheorem{theorem}{Theorem}
\newtheorem{lemma}{Lemma}
\newcommand{\C}{\mathbb C}
\newcommand{\R}{\mathbb R}
\newcommand{\ord}{\operatorname{ord}}
\newcommand{\ceil}[1]{\left\lceil #1\right\rceil}
\newcommand{\floor}[1]{\left\lfloor #1\right\rfloor}

\title{A Counterexample to the Conjectured Multiplicity\\
at the Forced Points at Infinity}
\author{Run \texttt{fa\_banach\_001}, lane 8\\
\texttt{agent\_lane\_08} (GPT5.6)}
\date{27 August 2026}

\begin{document}
\maketitle

\begin{center}
\fbox{\parbox{0.91\linewidth}{
\textbf{Status:} \texttt{candidate\_full\_negative\_solution\_likely\_valid}.\\
The conjectured value is false.  The sharp generic multiplicity is
$\ceil{(d\bmod n)/2}$, and this smaller value occurs inside the strictly
hyperbolic locus for every residue covered by the conjecture.
}}
\end{center}

\begin{abstract}
Simon and Vinzant conjecture that, when $r=d\bmod n\geq3$, the two forced
complex singularities $[0:1:\pm i]$ of a $C_n$-invariant hyperbolic plane
curve each have multiplicity $\binom r2$.  We show that the invariant
monomial support forces only the sharp lower bound $\ceil{r/2}$ and that this
is the generic exact multiplicity.  For every $d=qn+r$ with $q\geq1$ and
$r\geq3$, we construct normalized strictly hyperbolic real forms, invariant
even under the associated dihedral symmetry, whose multiplicity at both
points is exactly $\ceil{r/2}$.  Since $\ceil{r/2}<\binom r2$, these are full
counterexamples.  In particular, for $(n,d)=(4,7)$ the multiplicity is $2$
rather than $3$.
\end{abstract}

\section{Source conjecture and evidence}

The source is Faye Pasley Simon and Cynthia Vinzant, \emph{Invariant
hyperbolic curves: determinantal representations and applications to the
numerical range}, arXiv:2102.01726 (2021) \cite{SimonVinzant}.  In Section 8.1,
PDF/printed page 20, the authors write that for $d\bmod n\geq3$ the invariant
curves have complex singularities at $[t:x:y]=[0:1:\pm i]$ and conjecture
that each has multiplicity
\[
    \binom{d\bmod n}{2}.
\]
They immediately recall the invariant monomial basis that we use below.  The
following is a direct crop rendered from that source page.

\begin{center}
\includegraphics[width=\linewidth]{figures/open_problem_crop.png}
\end{center}

\paragraph{Precise transcription.}
``For curves with $d\bmod n\geq3$, it seems there are always multiple
singularities at infinity [\ldots]. More specifically, there are complex
singularities at the points $[t:x:y]=[0:1:\pm i]$. We conjecture they each
have multiplicity $\binom{d\bmod n}{2}$.''

\section{Definitions and local coordinates}

Put
\[
    u=x+iy,\qquad v=x-iy.
\]
The degree-$d$ invariant space has the basis recorded in the source:
\begin{equation}\label{eq:basis}
 \C[t,x,y]^{C_n}_d
 =\operatorname{span}_{\C}
 \bigl\{t^{d-j-k}u^jv^k:
 j\equiv k\pmod n,\ j+k\leq d\bigr\}.
\end{equation}
For a homogeneous plane curve $f=0$, its multiplicity at a projective point
is the least total degree of a nonzero term in a dehomogenized local Taylor
expansion.

At $P_+=[0:1:i]$, work in the chart $x=1$ and write $t=\tau$, $y=i+z$.
Then
\[
   u=iz,\qquad v=2-iz,
\]
so $v$ is a local unit.  Consequently
\begin{equation}\label{eq:localorder}
 \ord_{P_+}\bigl(t^{d-j-k}u^jv^k\bigr)
 =(d-j-k)+j=d-k.
\end{equation}
At $P_-=[0:1:-i]$, the roles of $u$ and $v$ are reversed and the order is
$d-j$.

\section{Proof intuition}

The source's own exponent diagram contains the obstruction to the proposed
number.  At $P_+$, maximizing the exponent $k$ minimizes local order by
\eqref{eq:localorder}.  The congruence $j\equiv k\pmod n$ and the degree
constraint allow $k$ to reach
$qn+\floor{r/2}$, not merely a value compatible with order $\binom r2$.
The corresponding support point is unique, so its coefficient cannot cancel
against another monomial.  This gives the exact generic order $\ceil{r/2}$.

To ensure this is not an artifact of leaving the hyperbolic locus, start with
a radial form having $d$ uniformly separated real roots in the $t$-direction
and add the two conjugate extremal monomials.  A compact sign argument keeps
all $d$ roots real and distinct for small perturbations, while the extremal
term lowers the local multiplicity to $\ceil{r/2}$ at both forced points.

\section{The sharp support calculation}

Write throughout
\[
   d=qn+r,\qquad 0\leq r<n,
   \qquad s=\floor{r/2},\qquad \delta=r-2s\in\{0,1\}.
\]

\begin{lemma}[extremal invariant monomial]\label{lem:support}
Among the pairs $(j,k)$ in \eqref{eq:basis},
\[
 \max k=qn+s,
\]
and equality holds only for $(j,k)=(s,qn+s)$.  Symmetrically,
$\max j=qn+s$, uniquely at $(j,k)=(qn+s,s)$.
\end{lemma}

\begin{proof}
Write $k=an+b$ with $0\leq b<n$.  Since $j\equiv k\pmod n$ and $j\geq0$,
we have $j\geq b$.  Hence
\[
 d\geq j+k\geq an+2b.
\]
In particular $a\leq q$.  If $a\leq q-1$, then
$k\leq(q-1)n+(n-1)=qn-1\leq qn+s$.  If $a=q$, the displayed inequality
gives $2b\leq r$, so $b\leq s$ and again $k\leq qn+s$.

The pair $(j,k)=(s,qn+s)$ is admissible because its exponent difference is
$qn$ and its sum is $qn+2s\leq qn+r=d$.  If $k=qn+s$, every nonnegative
$j\equiv s\pmod n$ has the form $s+\ell n$; the degree constraint gives
$\ell n\leq r-2s=\delta<n$, hence $\ell=0$.  This proves uniqueness.  The
statement for $j$ follows by symmetry.
\end{proof}

\begin{theorem}[sharp forced and generic multiplicity]\label{thm:generic}
Every $f\in\C[t,x,y]^{C_n}_d$ vanishes at $P_+$ and $P_-$ to order at least
$\ceil{r/2}$ whenever $r>0$.  This bound is sharp.  More precisely, the
coefficient of the unique extremal monomial
\[
 t^\delta u^s v^{qn+s}
\]
determines the lowest homogeneous part at $P_+$, and the coefficient of its
swapped monomial determines it at $P_-$.  Thus a Zariski-open dense set of
invariant forms has exact multiplicity $\ceil{r/2}$ at both points.
\end{theorem}

\begin{proof}
By \eqref{eq:localorder} and Lemma~\ref{lem:support}, every invariant basis
monomial has order at least
\[
 d-(qn+s)=r-s=\ceil{r/2}
\]
at $P_+$.  Equality comes from the unique stated monomial.  Therefore its
nonzero coefficient gives exact order and cannot cancel.  Swapping $j,k$
proves the assertion at $P_-$.  Nonvanishing of each of the two extremal
coefficients is the complement of a pair of coefficient hyperplanes, hence
is Zariski open and dense.
\end{proof}

\section{Strictly hyperbolic counterexamples}

Recall that a real homogeneous degree-$d$ form $f$ is strictly hyperbolic
with respect to $e=(1,0,0)$ if $f(e)>0$ and, for every nonzero $(x,y)$, the
polynomial $t\mapsto f(t,x,y)$ has $d$ distinct real roots.

\begin{theorem}[counterexamples for every conjectured residue]\label{thm:counter}
Let $d=qn+r$ with $q\geq1$ and $3\leq r<n$.  There exists $\varepsilon_0>0$
such that, for every real $\varepsilon$ with
$0<|\varepsilon|<\varepsilon_0$, the form $F_\varepsilon$ defined below is
real, normalized by $F_\varepsilon(1,0,0)=1$, strictly hyperbolic, and
$C_n$-invariant (indeed invariant under swapping $u,v$), while
\[
 \operatorname{mult}_{P_+}(F_\varepsilon)
 =\operatorname{mult}_{P_-}(F_\varepsilon)
 =\ceil{r/2}<\binom r2.
\]
Consequently the multiplicity conjecture in Section 8.1 of
\cite{SimonVinzant} is false.
\end{theorem}

\begin{proof}
Let $m=\floor{d/2}$ and define the radial form
\[
 H_d(t,x,y)=
 \begin{cases}
 \displaystyle\prod_{a=1}^{m}\bigl(t^2-a^2uv\bigr),&d=2m,\\[6pt]
 \displaystyle t\prod_{a=1}^{m}\bigl(t^2-a^2uv\bigr),&d=2m+1.
 \end{cases}
\]
Define the extremal real invariant form
\[
 G_{n,d}=t^\delta
 \left(u^s v^{qn+s}+u^{qn+s}v^s\right),
 \qquad F_\varepsilon=H_d+\varepsilon G_{n,d}.
\]
Both summands have degree $d$.  Complex conjugation swaps $u,v$, so
$G_{n,d}$ has real coefficients.  Under the generator of $C_n$, $u$ and $v$
acquire reciprocal $n$th-root factors; the two exponent differences are
$\pm qn$, so $G_{n,d}$ is invariant.  It is also fixed by the reflection
$u\leftrightarrow v$.  The same properties are immediate for $H_d$, and
$F_\varepsilon(1,0,0)=1$.

At $P_+$, the first monomial of $G_{n,d}$ has local order
$\delta+s=\ceil{r/2}$ and a nonzero initial term, while the second has larger
order.  At $P_-$ their roles reverse.  Each factor
$t^2-a^2uv$ has local order $1$ at either point, so $H_d$ has order
$\ceil{d/2}$.  Since $q\geq1$ and $r<n$, we have
$\ceil{d/2}>\ceil{r/2}$.  Thus every $\varepsilon\ne0$ gives exact
multiplicity $\ceil{r/2}$ at both points.

It remains to preserve strict hyperbolicity.  On the unit circle
$x^2+y^2=uv=1$, $H_d(t,x,y)$ is the fixed monic polynomial
\[
 Q_d(t)=
 \begin{cases}
 \prod_{a=1}^{m}(t^2-a^2),&d=2m,\\
 t\prod_{a=1}^{m}(t^2-a^2),&d=2m+1,
 \end{cases}
\]
whose $d$ roots are distinct and real.  Choose pairwise disjoint small closed
intervals around those roots such that $Q_d$ has opposite nonzero signs at
the two endpoints of every interval.  There are finitely many endpoints,
and $G_{n,d}$ is continuous on their product with the unit circle.  Hence,
for all sufficiently small $|\varepsilon|$, the perturbation does not change
any endpoint sign.  The intermediate value theorem then supplies a real root
of $F_\varepsilon(\,cdot\,,x,y)$ in each of the $d$ disjoint intervals.
Because the polynomial is monic of degree $d$, these are all its roots and
are distinct.  Homogeneity reduces every nonzero $(x,y)$ to the unit circle,
proving strict hyperbolicity.

Finally, $\ceil{r/2}<r(r-1)/2=\binom r2$ for every $r\geq3$, so the exact
multiplicity contradicts the conjectured value.
\end{proof}

\subsection*{The concrete $(n,d)=(4,7)$ example}

Here $q=1$, $r=3$, $s=1$, and $\delta=1$.  The construction is
\[
 H=t(t^2-uv)(t^2-4uv)(t^2-9uv),\qquad
 G=t(uv^5+u^5v),\qquad F_\varepsilon=H+\varepsilon G.
\]
At $P_+$, using $u=iz$, $v=2-iz$, the initial part of $G$ is
$32i\,\tau z$, of total degree $2$, whereas $H$ begins in degree $4$.
At $P_-$ the initial part is $-32i\,\tau z$.  Thus both multiplicities are
$2$, not $\binom32=3$.

\section{Verification, novelty, and limitations}

\paragraph{Exact checks.}
The included script \texttt{code/verify\_local\_orders.py} expands the
concrete example at both points, checks orders $4$ and $2$, verifies real
coefficients and quarter-turn/reflection invariance, and brute-forces the
support formula over a finite grid.  These checks are not used as proof; the
general argument is Lemma~\ref{lem:support} and Theorems~\ref{thm:generic}--
\ref{thm:counter}.  Full commands and outputs are recorded in
\texttt{verification.md}.

\paragraph{Bounded novelty search.}
On 27 August 2026, the run's registry, solution, attempt, and proof-gap
indexes were searched by arXiv id and core terms.  Live web/arXiv searches
used the exact conjecture sentence, title, authors, arXiv id, and combinations
of ``$d$ mod $n$'', ``multiplicity'', ``singularity'', ``correction'', and
``counterexample''.  They found the source and mirrors, but no later paper or
correction resolving this conjecture.  This is a bounded search, not an
exhaustive novelty guarantee.

\paragraph{Scope and interpretation.}
The argument uses ordinary plane-curve multiplicity, namely local Taylor
order at the named projective points.  It disproves both a universal reading
and a generic reading: Theorem~\ref{thm:generic} gives the generic value, and
Theorem~\ref{thm:counter} places it inside the strict hyperbolicity locus.
No assertion is made here about the distinct determinantal-representation
question that follows later in Section 8.1 of the source.

\paragraph{Human review recommendation.}
Check equation \eqref{eq:localorder}, the uniqueness clause in
Lemma~\ref{lem:support}, and the compact endpoint-sign argument for strict
hyperbolicity.  No external result beyond the source's invariant basis is
needed.

\begin{thebibliography}{9}
\bibitem{SimonVinzant}
F. Pasley Simon and C. Vinzant,
\emph{Invariant hyperbolic curves: determinantal representations and
applications to the numerical range}, arXiv:2102.01726, 2021.
\end{thebibliography}

\end{document}
